\chapter{Antes de começar}
\label{cap:antes}

\section{Por que versionar?}

Quem já trabalhou num documento importante conhece esta pasta:

\begin{center}
\begin{tikzpicture}[doc/.style={draw=azulMedio, fill=cinzaSuave, rounded corners=2pt,
    font=\ttfamily\small, anchor=west, minimum width=9cm, inner sep=5pt}]
  \node[doc] at (0,0) {\faFile*[regular]\; relatorio.docx};
  \node[doc] at (0,-0.75) {\faFile*[regular]\; relatorio\_v2.docx};
  \node[doc] at (0,-1.5) {\faFile*[regular]\; relatorio\_v2\_final.docx};
  \node[doc] at (0,-2.25) {\faFile*[regular]\; relatorio\_v2\_final\_AGORA\_VAI.docx};
  \node[doc, draw=vermelhoAlerta, fill=vermelhoClaro] at (0,-3) {\faFile*[regular]\; relatorio\_v2\_final\_corrigido(1).docx};
\end{tikzpicture}
\end{center}

Ninguém sabe qual é a versão certa, quem mudou o quê, nem por quê. Agora imagine isso com um sistema de milhares de arquivos e várias pessoas mexendo ao mesmo tempo. É esse o problema que um \termo{sistema de controle de versão} resolve.

Com controle de versão, cada alteração fica registrada com:

\begin{itemize}
  \item \textbf{o que} mudou, linha a linha;
  \item \textbf{quem} mudou e \textbf{quando};
  \item \textbf{por que} mudou, numa mensagem escrita por quem fez;
  \item a possibilidade de \textbf{voltar} a qualquer estado anterior;
  \item várias pessoas trabalhando \textbf{em paralelo} sem apagar o trabalho umas das outras.
\end{itemize}

\begin{analogia}
Pense no Git como o \enquote{histórico de versões} de um documento compartilhado, só que muito mais poderoso: ele guarda o projeto inteiro, funciona sem internet e deixa cada pessoa trabalhar numa cópia própria antes de juntar tudo.
\end{analogia}

\section{O terminal, sem medo}

Quase tudo nesta apostila é feito no \termo{terminal} (também chamado de linha de comando, console ou \emph{shell}). É uma janela onde você digita uma ordem e aperta \tecla{Enter}.

\begin{itemize}
  \item \textbf{Linux}: procure por \enquote{Terminal} no menu de aplicativos.
  \item \textbf{Windows}: depois de instalar o Git, use o \textbf{Git Bash}, que vem junto. Ele aceita os mesmos comandos mostrados aqui.
\end{itemize}

Cinco comandos básicos ajudam a se localizar:

\begin{center}
\renewcommand{\arraystretch}{1.3}
\begin{tabularx}{0.9\linewidth}{@{}>{\ttfamily\color{azulMB}}l X@{}}
\toprule
pwd & mostra em que pasta você está \\
ls & lista os arquivos da pasta atual \\
cd pasta & entra numa pasta (\texttt{cd ..} volta uma pasta) \\
mkdir nome & cria uma pasta \\
cat arquivo & mostra o conteúdo de um arquivo \\
\bottomrule
\end{tabularx}
\end{center}

\section{Instalando o Git}

\begin{terminalt}{· LINUX (DEBIAN/UBUNTU)}
sudo apt update
sudo apt install git
\end{terminalt}

No \textbf{Windows}, baixe o instalador oficial em \url{https://git-scm.com/downloads} e aceite as opções padrão. Ao final, abra o \textbf{Git Bash}.

Para conferir se deu certo:

\begin{terminal}
git --version
\end{terminal}
\begin{saida}
git version 2.34.1
\end{saida}

Qualquer versão a partir da 2.28 funciona com todos os exemplos desta apostila. Se a sua for mais antiga, veja a dica da seção \ref{sec:primeiro-repo}.

\section{Dizendo ao Git quem é você}

Todo registro no Git leva o nome e o e-mail de quem o fez. Configure uma única vez por computador:

\begin{terminal}
git config --global user.name "Seu Nome"
git config --global user.email "seu.email@exemplo.com"
# faz a branch inicial de novos repositórios se chamar main
git config --global init.defaultBranch main
\end{terminal}

\begin{dica}
Use o mesmo e-mail cadastrado no GitLab da instituição. Assim o GitLab reconhece os seus commits e mostra a sua foto e o seu nome ao lado deles.
\end{dica}

Para ver o que está configurado:

\begin{terminal}
git config --list
\end{terminal}
